\chapter{Robert K. Merton, G. H. Mead \& the Sociology of Religion}

\section{Merton --- Role, Status \& Critique of Structural Functionalism}

\subsection{Answer-Writing: European vs American Sociology}

The lecture opens with a 10-mark (2013) question: \textbf{``Sociology emerged in Europe and flourished, to begin with, on a social reformist orientation in the USA.'' (150 words)}

\begin{PYQLink}
\begin{itemize}
\item This is a 2013 UPSC question (10 marks). It is essentially asking for the \textbf{European vs American sociology} contrast.
\end{itemize}
\end{PYQLink}

\begin{ComparisonBox}
\begin{itemize}
\item \textbf{European sociology} emerged in response to \textbf{intellectual changes} (Renaissance, Enlightenment, Scientific Revolution, French Revolution, Industrial Revolution --- social change over $\sim$400 years, 1300--1700). It was \textbf{large-scale, theoretical} (grand theory) --- Durkheim on collective consciousness and weakening cohesion, Marx, Weber --- and moved from \textbf{macro to micro}.
\item \textbf{American sociology} emerged in response to \textbf{social issues} (migration, crime rate, homelessness, poverty) within the ``golden triangle'' (Chicago--Boston--Washington). It was \textbf{small-scale, empirical, pragmatic, field-study-based and interactionist} (verstehen --- understanding issues from people's perspective to \textit{address} them, not just study them). It started \textbf{micro and moved to macro}, then to the \textbf{middle range} (Merton).
\item In Europe, \textbf{interpretivism was a challenge to the collectivist orientation}; in America, \textbf{collectivism was the challenge to interpretivism} --- the intellectual challenge ran in opposite directions. (Both \textbf{Parsons} and \textbf{Weber} moved from \textit{social action} to \textit{social structures} --- Parsons from the structure of social action to the social system; Weber from social action to large-scale structures like capitalism.)
\end{itemize}
\end{ComparisonBox}

\begin{CriticalPoint}
\begin{itemize}
\item Marx said the point of philosophy so far has been to \textit{interpret} the world, but it should be to \textit{change} it --- the idea of \textbf{activism} comes from Marx (an action theorist). Yet he analysed structurally (``communism is structurally inevitable''); it was \textbf{Lenin} who gave spirit to Marx's writings --- revolution through a political party, the \textbf{vanguard party}.
\end{itemize}
\end{CriticalPoint}

\begin{DataFact}
\begin{itemize}
\item Indian sociology emerged in the context of \textbf{colonialism} --- a very different context again.
\end{itemize}
\end{DataFact}

\begin{QuickRevision}
\begin{itemize}
\item Europe = intellectual changes $\rightarrow$ grand, theoretical, macro$\rightarrow$micro. USA = social issues $\rightarrow$ empirical, pragmatic, micro$\rightarrow$macro$\rightarrow$middle range.
\item Marx = action/activism (change the world), but structural analyst; Lenin = vanguard party.
\item India = colonialism context.
\end{itemize}
\end{QuickRevision}

\sectiondivider

\subsection{Role \& Status: Role Set and Status Set}

\begin{KeyConcept}
\begin{itemize}
\item The concept of role and status was discussed by \textbf{psychologists} and \textbf{anthropologists} before moving into sociology, and it \textbf{bridges sociology and psychology}. Scholars: \textbf{Ralph Linton}, \textbf{Parsons}, \textbf{Luhmann}, \textbf{Merton}.
\item Start with two basic concepts: \textbf{role set} and \textbf{status set}.
\end{itemize}
\end{KeyConcept}

\begin{KeyConcept}
\textbf{Status set:}
\begin{itemize}
\item \textbf{Status} = \textit{who we are} (what we are in society). I am a teacher (achieved) and a son (ascribed) at the same time; I hold \textbf{both} ascribed and achieved statuses.
\item \textbf{Ascribed status} = status you are born with (son, brother, friend, colleague). \textbf{Achieved status} = status you achieve (engineer, teacher, sociologist, researcher).
\item \textbf{Status set} = the set of \textbf{ascribed + achieved} statuses a person holds at a given time.
\end{itemize}
\end{KeyConcept}

\begin{KeyConcept}
\textbf{Role set:}
\begin{itemize}
\item \textbf{Role} = \textit{what we do}. Every status has \textbf{multiple roles}.
\item As a \textbf{teacher} (status), I have many roles: conduct class, conduct tests, organise notes, help with daily answer-writing, run test series/courses, resolve doubts (doubt sessions). As a \textbf{university teacher}, I also publish journals and write articles (which add to my profile for promotion).
\item As a \textbf{mother} (status), roles include: care for husband/children, cook, attend parent-teacher meetings, manage the household budget, hire help for cleaning/dish-washing.
\item \textbf{Role set} = all the roles linked to a status.
\end{itemize}
\end{KeyConcept}

\begin{CriticalPoint}
\begin{itemize}
\item Some statuses are neither clearly ascribed nor achieved. \textbf{Motherhood} is actually achieved (you are not a mother by birth), yet we call it ascribed --- this is a criticism of the binary.
\end{itemize}
\end{CriticalPoint}

\begin{QuickRevision}
\begin{itemize}
\item Status = what we are; Role = what we do.
\item \textbf{Status set} = ascribed + achieved statuses. \textbf{Role set} = roles attached to one status.
\item Role bridges sociology and psychology; scholars Linton, Parsons, Luhmann, Merton.
\end{itemize}
\end{QuickRevision}

\sectiondivider

\subsection{Role Conflict}

\begin{KeyConcept}
\textbf{Role conflict:}
\begin{itemize}
\item \textbf{Role conflict} arises when the roles of two different statuses clash. A mother who is an IPS officer must, at home, be a mother (particular, emotional) but in society be an IPS officer (universal, \textbf{affective-neutral}, instrumental --- pattern variable B standards).
\end{itemize}
\end{KeyConcept}

\begin{ExampleBox}
\textbf{The IPS-officer mother:}
\begin{itemize}
\item A mother is also a teacher; her son happens to be her student and is failing. As a mother she wants to help her son; as a teacher she must maintain \textbf{universal, honest} standards and the image of impartiality. She must balance both --- this is \textbf{role conflict} (she finally passes him at 33\%, balancing her motherly concern and her professional image).
\end{itemize}
\end{ExampleBox}

\begin{ComparisonBox}
\textbf{The diagram (status set + role set):}
\begin{itemize}
\item Inner circle = \textbf{status set} (professor, mother, researcher, wife). Outer edge = \textbf{role sets}.
\item Professor $\rightarrow$ teacher role + colleague role (discuss, attend seminars/conferences). Mother $\rightarrow$ maternal role + civic role. Researcher $\rightarrow$ author role + field-worker role. Wife $\rightarrow$ marital role + domestic role.
\end{itemize}
\end{ComparisonBox}

\begin{ExampleBox}
\textbf{Relevance to reference group theory:}
\begin{itemize}
\item A mother who wants to become a teacher has ``teacher'' as her \textbf{reference group}. Through \textbf{anticipatory socialisation} she joins it; but if she favours her son, other teachers complain to the principal and she is suspended --- she is reduced to \textbf{marginality} (Robert Ezra Park): neither a proper mother nor a proper teacher in that context.
\end{itemize}
\end{ExampleBox}

\begin{QuickRevision}
\begin{itemize}
\item \textbf{Role conflict} = clash between roles of different statuses (e.g., IPS officer who is also a mother; teacher whose son is her student).
\item Status set (inner) + role sets (outer); reference group $\rightarrow$ anticipatory socialisation $\rightarrow$ marginality if roles clash.
\end{itemize}
\end{QuickRevision}

\sectiondivider

\subsection{Merton's Critique of Structural Functionalism}

\begin{CriticalPoint}
\textbf{Who is being criticised?}
\begin{itemize}
\item This is \textbf{not} Merton criticising Talcott Parsons. Some elements may look like a critique of Parsons, but largely Merton developed the critique of structural functionalism to criticise \textbf{Malinowski and Radcliffe-Brown} (the anthropological structural functionalists), not Parsons.
\end{itemize}
\end{CriticalPoint}

\begin{DataFact}
\textbf{The origin of structural functionalism:}
\begin{itemize}
\item Structural functionalism as a theory originates in \textbf{Herbert Spencer} (who never called it that, but was describing it). It was later used by \textbf{Durkheim}, who inspired \textbf{Radcliffe-Brown}. The two dominant names are \textbf{Durkheim and Parsons}; Radcliffe-Brown and Durkheim were contemporaries.
\item \textbf{Anthropologist vs sociologist} --- conventionally anthropologists study simple/tribal societies, sociologists study modern/urban societies. But sociologists also study tribes and villages, and anthropologists study urbanisation and nationalism. The \textbf{only difference is degree/emphasis} --- sociologists emphasise modernity, anthropologists emphasise simple societies. (Anthropologists trace today's practices to the ancient past --- e.g., blaming Muslims during COVID = ``witch-hunting'' = what \textbf{Stanley Cohen} calls \textbf{folk devilization}.)
\end{itemize}
\end{DataFact}

\begin{QuickRevision}
\begin{itemize}
\item Merton critiques the \textit{anthropological} structural functionalism (Malinowski, Radcliffe-Brown), not Parsons.
\item Origin: Herbert Spencer $\rightarrow$ Durkheim $\rightarrow$ Radcliffe-Brown.
\item Anthropologist vs sociologist = only a difference of degree/emphasis.
\end{itemize}
\end{QuickRevision}

\sectiondivider

\subsection{The Three Functional Postulates}

Merton found issues with \textbf{three basic postulates} of (anthropological) structural functionalism:

\begin{DataFact}
\textbf{The three postulates (memorise):}
\begin{enumerate}
\item \textbf{Functional unity postulate} --- any given social institution/structure plays a role in the \textit{maintenance of the system}.
\item \textbf{Functional indispensability postulate} --- a given social institution/structure has \textit{no alternative/no substitute} and \textit{has to be there} to perform its requisite function.
\item \textbf{Functional universality postulate} --- every social institution/structure plays a \textit{positive} function.
\end{enumerate}
\end{DataFact}

\paragraph{Merton's critique of each postulate}

\begin{CriticalPoint}
\textbf{Critique of functional universality:}
\begin{itemize}
\item Religion (and the state) do not play only positive functions. Religion also plays \textbf{negative} functions: communal politics, fundamentalism, terrorism, secessionist movements, sectarian wars. The state too has negative functions: surveillance/snooping (e.g., Pegasus), raising fuel prices, interstate conflict whose costs citizens bear.
\item Functionalists, by always seeing the positive, are regarded as \textbf{pro-state and status-quoist}. Merton writes negative functions as \textbf{dysfunctions}.
\item Merton is thereby \textbf{bringing Marxism into structural functionalism}: structural functionalism read only the positive, Marxism only the negative --- but reality is \textbf{both}, so study both.
\end{itemize}
\end{CriticalPoint}

\begin{CriticalPoint}
\textbf{Critique of functional indispensability:}
\begin{itemize}
\item Nothing is indispensable; many institutions can do what one institution does. Punishment is given by the \textbf{judiciary}, but also by \textbf{religion} (blasphemy), by the \textbf{mob} (lynching), and by the \textbf{family} (slap, expulsion). (We do not judge which is right --- a sociologist simply records that punishment is given by many institutions.)
\item Marriage: the family's role can be done by \textbf{R.A. Samaj / Arya Samaj}, by the \textbf{court}, by \textbf{apps}; and the \textbf{state} decides who may marry whom (child marriage is void; the UP government's marriage law). Marx: \textbf{communism} will perform the roles religion once did, fairly.
\item \textbf{Judicial activism} = judiciary playing parliament's role; the \textbf{Ninth Amendment} = parliament playing the judiciary's role. No institution is indispensable.
\item Here Merton \textbf{shakes hands with Parsons} (structural differentiation already showed this) --- so it is wrong to call this a critique of Parsons.
\end{itemize}
\end{CriticalPoint}

\begin{CriticalPoint}
\textbf{Critique of functional unity:}
\begin{itemize}
\item Not every role played by an institution maintains the system; some \textbf{disintegrate} it. In \textit{simple} societies one religion (animism) unifies; in \textit{complex} societies multiple religions coexist and may disintegrate --- religious wars, the \textbf{Charlie Hebdo} case, the \textbf{Partition} of India--Pakistan, the Israel--Palestine conflict.
\item So the functional-unity postulate holds only for simple/tribal societies, not for modern complex societies.
\end{itemize}
\end{CriticalPoint}

\begin{CriticalPoint}
\textbf{The ``scientists are innovators'' koochi error:}
\begin{itemize}
\item Merton's ``innovator'' is not about scientists; a coaching guide (``koochi'') wrongly lists scientists as innovators. Burn such guides (as Ambedkar said burn the shastras).
\end{itemize}
\end{CriticalPoint}

\begin{QuickRevision}
\begin{itemize}
\item Three postulates: \textbf{functional unity} (all maintain the system), \textbf{indispensability} (no substitute), \textbf{universality} (all positive).
\item Merton: universality is false (there are \textbf{dysfunctions}); nothing is indispensable (punishment/marriage by many institutions; judicial activism); unity holds only in simple societies (complex societies disintegrate).
\end{itemize}
\end{QuickRevision}

\sectiondivider

\subsection{Merton's Advisory: Net Balance Analysis \& Levels of Analysis}

\begin{KeyConcept}
\textbf{Merton's advisory (how a sociologist should analyse):}
\begin{itemize}
\item Statements like ``this is functional to the system'' are \textbf{vague, arbitrary and abstract}. What is functional to the system may be dysfunctional to some individuals; what is dysfunctional to some may be functional to others; and non-functional to yet others.
\end{itemize}
\end{KeyConcept}

\begin{ExampleBox}
\textbf{Slavery:}
\begin{itemize}
\item Slavery was functional to the \textbf{system} (productivity), but dysfunctional to the \textbf{slaves} (and only truly functional to the masters). So we must ask: functional to \textit{whom}?
\end{itemize}
\end{ExampleBox}

\begin{KeyConcept}
\textbf{Net balance analysis:}
\begin{itemize}
\item Always conduct a \textbf{net balance analysis}: study functions, but ask \textbf{functional to whom?} (system, group, individuals, individuals-within-other-groups). Example: \textbf{unpaid domestic work} --- functional to male members/family, dysfunctional to the women doing it (and differently for North vs South Indian, Western vs Indian women).
\end{itemize}
\end{KeyConcept}

\begin{KeyConcept}
\textbf{Levels of analysis:}
\begin{itemize}
\item Develop \textbf{levels of analysis} (system / group / individual). Anything \textbf{patterned, regular and standardised} can be subjected to structural-functional analysis --- not just large institutions. Merton thereby \textit{liberated} structural functionalism to study empirical social processes, cultural forms, norms and values.
\end{itemize}
\end{KeyConcept}

\begin{ExampleBox}
\textbf{Durga Puja, poverty, nationalism, cricket:}
\begin{itemize}
\item \textbf{Durga Puja} is functional to Bengalis/Hindus, dysfunctional to car drivers stuck in traffic, and non-functional to America/Jains. \item \textbf{Poverty} is functional to politicians (vote bank, development discourse) and industrialists (cheap labour = industrial reserve army), but dysfunctional to the poor and to India's image abroad (while attracting FDI as cheap-labour source). \item \textbf{Nationalism} may be functional to one group and dysfunctional to another (nuclear arms race; Tagore called nationalism an ``evil epidemic''). \item A \textbf{cricket match} is functional to bookies, TV/advertising agencies, and stressed students, but dysfunctional to other sports (hockey) and to the spirit of sportsmanship (creating celebrities, not sportsmen).
\item There is \textbf{no objective reality} --- it depends on the perspective (``which glasses you see the world with'').
\end{itemize}
\end{ExampleBox}

\begin{CriticalPoint}
\begin{itemize}
\item When reading Merton, never claim anything is functional \textit{to the system as a whole} --- it is always functional \textit{to groups}. (Merton is a scholar of groups --- reference groups, relative deprivation.)
\end{itemize}
\end{CriticalPoint}

\begin{QuickRevision}
\begin{itemize}
\item Statements about ``function'' are vague $\rightarrow$ conduct \textbf{net balance analysis} (functional to whom? dysfunctional to whom? non-functional to whom?) with \textbf{levels of analysis}.
\item Anything patterned/regular/standardised can be studied (not just institutions).
\item No objective reality; Merton studies groups, not ``the system.''
\end{itemize}
\end{QuickRevision}

\sectiondivider

\subsection{Manifest and Latent Functions}

\begin{KeyConcept}
\begin{itemize}
\item \textbf{Manifest} = visible/anticipated. \textbf{Latent} = hidden/unanticipated. The \textbf{task of the sociologist is to reveal what is hidden} (this is what makes sociology a science, as \textbf{Pierre Bourdieu} says --- science reveals what is hidden).
\item \textbf{Manifest function} = the \textit{anticipated} consequence. \textbf{Latent function} = the \textit{unanticipated} consequence. Unanticipated consequence has \textbf{three forms}: latent function, \textbf{dysfunction}, and \textbf{non-function} (the latter is less important to Merton).
\item \textbf{Colin Campbell's critique:} ``latent'' is a little vague --- what is hidden to one person may be anticipated/obvious to another (latent is subjective).
\end{itemize}
\end{KeyConcept}

\begin{ExampleBox}
\textbf{The examination (manifest vs latent):}
\begin{itemize}
\item \textbf{Manifest functions} (anticipated): test knowledge, impart skills, find the best/meritorious, release merit lists, give jobs. (This is common sense.)
\item \textbf{Latent functions} (hidden): converts \textbf{social privilege into merit} (examination is a machine that turns privilege into merit --- hence the absence of SC/ST students and faculty in IITs/IIMs); prepares the \textbf{marriage market} (passing an exam like IAS/IES makes one a good match); \textbf{hides unemployment} (it ``curtains'' unemployment by keeping people preparing).
\item \textbf{Dysfunctions}: promotes \textbf{business out of students' fear} and proliferates coaching centres; the rise of criminality is a dysfunction (those without legitimate means/education turn to crime --- anomie is rooted in social structure); the \textbf{digital divide}; producing rote learners (Gandhi and Tagore said this); establishing an authority relation between examiner and examined.
\end{itemize}
\end{ExampleBox}

\begin{CriticalPoint}
\textbf{Subject matter of sociology:}
\begin{itemize}
\item The subject matter of sociology, per Merton, is the study of the \textbf{latent}. (Weber = social action; Durkheim = social fact; Marx = class struggle; Parsons = social system; Merton = latent.)
\item Merton found the common ``spirit of sociology'' (what C. Wright Mills called \textbf{sociological imagination}) across Marx, Weber, Durkheim and Parsons --- all studied the hidden (Weber's meanings behind work; Freud's unconscious; Durkheim's collective effervescence). He learned this from Weber but made it explicit.
\item \textbf{Robert King Merton} --- ``King'' because he was once a \textbf{street magician} (magicians wear a king's crown), and a magician \textit{reveals the hidden}.
\end{itemize}
\end{CriticalPoint}

\begin{ExampleBox}
\textbf{Social distancing (net empirical analysis):}
\begin{itemize}
\item Social distancing is now \textbf{patterned/regularised} and can be studied. It is functional to everyone (avoid death), but dysfunctional to universities (classes suspended, board exams suspended for the first time, curriculum changed) and to personality-test candidates; it may even be functional to some candidates (more preparation time). So do not say ``functional'' or ``dysfunctional'' \textit{only} --- that is value judgment; do a \textbf{net empirical analysis}.
\end{itemize}
\end{ExampleBox}

\begin{QuickRevision}
\begin{itemize}
\item Manifest = anticipated; Latent = unanticipated (latent function + dysfunction + non-function).
\item Sociology's subject matter = the \textbf{latent}; sociologist reveals the hidden.
\item Examination: manifest (test knowledge, jobs) vs latent (privilege$\rightarrow$merit, marriage market, hides unemployment) + dysfunctions (coaching business, criminality, digital divide).
\end{itemize}
\end{QuickRevision}

\sectiondivider

\section{Merton's Critique of Deviance \& G. H. Mead}

\subsection{Answer-Writing: Fact vs Value in Weber's Protestant Ethic}

The opening question (10 marks): \textbf{``Distinguish between fact and value in Weber's \textit{The Protestant Ethic and the Spirit of Capitalism}.''}

\begin{PYQLink}
\begin{itemize}
\item 10-mark question; answer by extracting from understanding. Technique: define keywords, or describe briefly; it depends on the demand and the marks.
\end{itemize}
\end{PYQLink}

\begin{DataFact}
\textbf{Facts (unbiased, value-free, recorded):}
\begin{itemize}
\item The \textbf{Industrial Revolution happened in England}; the \textbf{Protestant Reformation happened in Germany}; \textbf{$\sim$70\% of the staff, workers, employees and managers during the industrial revolution were Protestants}. These are recorded, undisputed facts. The \textbf{schism} (division) of the Catholic church into Catholic and Protestant churches after the Reformation is also a fact.
\item These three facts can be summed up as the \textbf{material conditions}.
\end{itemize}
\end{DataFact}

\begin{KeyConcept}
\textbf{Where value enters:}
\begin{itemize}
\item \textit{Why} did these facts arise? The answer lies in \textbf{Puritan work ethics / Protestant ideas}. So the ideas (values) explain the facts, and the relation between fact and value is \textbf{adequate causal} (Weber) --- a multi-causal reality (rational laws, markets are also factors), but it is the \textit{Protestant values} that produced rational capitalism.
\end{itemize}
\end{KeyConcept}

\begin{DataFact}
\textbf{Coleman's boat (James Coleman):}
\begin{itemize}
\item Coleman depicted Weber's idea as a ``boat'': \textbf{rational capitalism} (fact) at the top, \textbf{Puritan work ethics} (values) at the bottom. Puritan work ethics $\rightarrow$ initially \textbf{value-rational social action} (work as spiritual experience, a way to heaven) $\rightarrow$ gradually became \textbf{zweck-rational} (instrumental) action (concerned with money, not God) $\rightarrow$ \textbf{formal rationality} $\rightarrow$ \textbf{rational capitalism}.
\end{itemize}
\end{DataFact}

\begin{ComparisonBox}
\textbf{Marx vs Weber on fact/value:}
\begin{itemize}
\item \textbf{Marx}: economic base = fact/material condition; values = superstructure; \textbf{monocausal} relation.
\item \textbf{Weber}: values are causal, facts are effects; relation is \textbf{probabilistic / adequate causal}, not monocausal.
\end{itemize}
\end{ComparisonBox}

\begin{QuickRevision}
\begin{itemize}
\item Facts (recorded): Industrial Revolution in England, Reformation in Germany, 70\% Protestants. Values: Protestant work ethic.
\item \textbf{Coleman's boat}: value-rational $\rightarrow$ zweck-rational (instrumental) $\rightarrow$ formal rationality $\rightarrow$ rational capitalism.
\item Marx (base$\rightarrow$superstructure, monocausal) vs Weber (values$\rightarrow$facts, adequate causal).
\end{itemize}
\end{QuickRevision}

\sectiondivider

\subsection{Merton: The Accidental Sociologist}

\begin{DataFact}
\textbf{Merton's biography:}
\begin{itemize}
\item Merton was an \textbf{accidental sociologist} --- a street magician/conjurer as a young man (his sister's boyfriend was a magician, his \textbf{role model / reference individual}). He got admission in \textbf{Temple University} (Philadelphia) as a poor, underprivileged person. There was no sociology discipline there.
\item He met \textbf{Pitirim Sorokin} (who had established Harvard's sociology department) through a teacher he fell in love with; Sorokin found potential in him and invited him to Harvard. Meeting Sorokin was \textbf{serendipity} (a happy accident).
\item By chance he attended a \textbf{Talcott Parsons} lecture; he could not understand it, but later Parsons' lectures came out as \textit{The Structure of Social Action} (which Merton could read easily because he had attended the lectures). He found flaws, reported them to Sorokin (his supervisor), and through Parsons got introduced to European sociology (Weber, Durkheim).
\item So Merton's life journey: an \textbf{innovator} (street magician) $\rightarrow$ \textbf{conformist} (student of sociology) $\rightarrow$ \textbf{rebellion} (criticising scholars and establishing his own name).
\end{itemize}
\end{DataFact}

\begin{DataFact}
\textbf{Merton's study of deviance:}
\begin{itemize}
\item Merton's study of deviance, like Durkheim's suicide, was based on \textbf{official statistics}. He found that the maximum deviants were from the \textbf{lower working class} (overpopulating American prisons); most were \textbf{blacks} --- so \textbf{class and race intertwined} to produce crime. Deviance, for Merton, is the result of the lower working class lacking legitimate means to gratify culturally defined goals.
\end{itemize}
\end{DataFact}

\begin{QuickRevision}
\begin{itemize}
\item Merton = accidental sociologist (magician $\rightarrow$ sociologist); serendipity with Sorokin; chance encounter with Parsons.
\item Deviance study = official statistics; lower working class + blacks; deviance = lack of legitimate means.
\end{itemize}
\end{QuickRevision}

\sectiondivider

\subsection{Critiques of Merton's Deviance Theory}

\begin{CriticalPoint}
\textbf{Albert Cohen (status frustration):}
\begin{itemize}
\item \textbf{Albert Cohen} (Merton's student) argued the lower working class turns to deviance not from lack of legitimate means but from \textbf{status frustration} --- however hard they work, they are not appreciated (the rickshaw-walla/auto-driver toils all day and gets only abuse; the Uber driver tries to escape the ``auto-walla'' image --- \textbf{Erving Goffman's role distance}).
\item This explains \textbf{vandalism} (crime for fun, e.g., by 30--50\% school ``failures'') which Merton's theory cannot explain.
\end{itemize}
\end{CriticalPoint}

\begin{CriticalPoint}
\textbf{Cloward \& Ohlin (illegitimate opportunity structure):}
\begin{itemize}
\item Merton described only the \textbf{legitimate opportunity structure} (socially legitimate means). But there is also an \textbf{illegitimate opportunity structure}: even in the world of crime there are goals and means, and one can fail there too. \textbf{Cloward and Ohlin} describe \textbf{three forms of criminal subculture}: \textbf{organised crime} (mafia --- with its own role models and competition), \textbf{conflict subculture} (turf wars, gang wars), and \textbf{retreatist subculture} (those who fail in both worlds $\rightarrow$ drug abuse).
\end{itemize}
\end{CriticalPoint}

\begin{CriticalPoint}
\textbf{Laurie Taylor:}
\begin{itemize}
\item \textbf{Laurie Taylor} asks: who defines the ``culturally defined goal''? The goal of material wealth is not \textit{culturally} defined but defined by the \textbf{capitalist class} (through Hollywood/media, creating aspirations and efficient workers for Facebook/Google/Microsoft). It is a trap, not a culture.
\item Also: many people lack legitimate means, yet only a few turn to deviance --- Merton does not explain this discrepancy.
\end{itemize}
\end{CriticalPoint}

\begin{QuickRevision}
\begin{itemize}
\item \textbf{Cohen}: status frustration (not lack of means) $\rightarrow$ vandalism.
\item \textbf{Cloward \& Ohlin}: illegitimate opportunity structure; three criminal subcultures (organised, conflict, retreatist).
\item \textbf{Laurie Taylor}: ``culturally defined goal'' is really the capitalist class's goal; and why only some deviate?
\end{itemize}
\end{QuickRevision}

\sectiondivider

\subsection{Urban vs Rural India: A Vindication of Merton}

\begin{DataFact}
\begin{itemize}
\item Urban India has a higher crime rate than rural India --- a fact that \textbf{vindicates Merton's theory} (anomie is rooted in social structure; urban India provides many material goals --- ``Mumbai, the city of dreams'', strugglers, e.g., the Raj Kundra case). America, more developed than India, also has a higher crime rate --- the more developed a society, the more pathological strains it may contain.
\item (Note: rural crime may go unreported, but urban under-reporting is equally or more a factor.)
\end{itemize}
\end{DataFact}

\begin{CriticalPoint}
\textbf{Bureaucratic personality:}
\begin{itemize}
\item Being an officer is not enough --- a \textbf{bureaucratic personality} is a specific type: rigid, inflexible about rules (means become ends; means-oriented, not goal-oriented) = the \textbf{ritualist}. \textbf{Rebellion} = those who bring social change.
\end{itemize}
\end{CriticalPoint}

\begin{QuickRevision}
\begin{itemize}
\item Urban India's higher crime rate vindicates Merton (more material goals $\rightarrow$ more anomie).
\item Bureaucratic personality = rigid/means-oriented = ritualist.
\end{itemize}
\end{QuickRevision}

\sectiondivider

\subsection{G. H. Mead: I, Me and Role-Taking}

\begin{DataFact}
\textbf{Who Mead is:}
\begin{itemize}
\item \textbf{George Herbert Mead} (1863--1931) is a \textbf{social psychologist} (not a sociologist), of the \textbf{Chicago school} (micro-sociology). He is the most affluent scholar in the syllabus (a gated community, wealthy family). Jane Addams, who got a Nobel Prize for activism, was also Chicago.
\end{itemize}
\end{DataFact}

\begin{KeyConcept}
\textbf{The basic question:}
\begin{itemize}
\item A two-year-old child takes a toffee happily; at 24 the same person refuses it --- because now she attaches \textit{meaning} to it (why is he giving me this?). Mead asks: \textbf{when did we start putting ourselves in the shoes of others and thinking about what they think of us?} This is \textbf{role-taking} (a term Mead coined): you take the role of the other, look back at yourself, and adjust the interaction.
\end{itemize}
\end{KeyConcept}

\begin{KeyConcept}
\textbf{I and Me:}
\begin{itemize}
\item \textbf{I} = the individual/impulsive self (the two-year-old who grabs the toffee). \textbf{Me} = the \textbf{social self} --- a highly \textbf{reflexive self} with the ability to take the other's position and change the course of interaction accordingly.
\end{itemize}
\end{KeyConcept}

\begin{CriticalPoint}
\textbf{Reflexivity:}
\begin{itemize}
\item \textbf{Reflexivity} (here) = the ability to take the other's position and ask ``why did he say this to me?'', then adjust. Mind-readers are the most reflexive. (Not everyone is reflexive.)
\end{itemize}
\end{CriticalPoint}

\begin{QuickRevision}
\begin{itemize}
\item Mead (1863--1931) = social psychologist, Chicago school, micro.
\item Basic question: when do we start role-taking (putting ourselves in others' shoes)?
\item \textbf{I} = impulsive self; \textbf{Me} = social/reflexive self.
\end{itemize}
\end{QuickRevision}

\sectiondivider

\subsection{Mead's Three Stages of the Social Self}

\begin{DataFact}
\textbf{The three stages (memorise):}
\begin{enumerate}
\item \textbf{Preparatory stage} (0--2 years) --- the actor performs roles \textit{without being conscious} of it (imitating).
\item \textbf{Play stage} (2--6 years) --- the child plays with the roles of \textbf{significant others} (mother, father) --- copying them (girl imitates mother cutting vegetables; boy imitates father).
\item \textbf{Game stage} (6 years onwards) --- the child develops the capacity to see himself from others' eyes and plays games using \textbf{role-taking of generalised others} (e.g., a batsman thinking where the bowler will pitch, or chess --- anticipating the other's move).
\end{enumerate}
\end{DataFact}

\begin{ExampleBox}
\textbf{The game stage:}
\begin{itemize}
\item In the play stage, putting on father's tie and being mocked (``you look funny'') makes the child \textbf{conscious} of how he looks to others. In the game stage, the batsman and bowler each put themselves in the other's shoes (the bowler pitches where the batsman might not expect; the batsman anticipates) --- role-taking with \textbf{generalised others}, not just significant others. This also creates the conflict between children and parents (the child now applies the same role-taking to parents).
\end{itemize}
\end{ExampleBox}

\begin{QuickRevision}
\begin{itemize}
\item \textbf{Preparatory} (0--2, unconscious imitation) $\rightarrow$ \textbf{Play} (2--6, significant others) $\rightarrow$ \textbf{Game} (6+, generalised others, role-taking).
\item Significant others = mother/father; generalised others = society.
\end{itemize}
\end{QuickRevision}

\sectiondivider

\subsection{Symbolic Interactionism \& Self and Society}

\begin{KeyConcept}
\textbf{Symbolic interactionism:}
\begin{itemize}
\item \textbf{Symbolic interactionism} was labelled by \textbf{Herbert Blumer}, not Mead. Every interaction is about reading \textbf{symbols} (the bowler's expression and speed; the bat's direction) and gauging the other's next move.
\end{itemize}
\end{KeyConcept}

\begin{KeyConcept}
\textbf{Self and society are twin-born:}
\begin{itemize}
\item Self and society are the same: society is made of selves, and the self develops in the context of society. \textbf{Self is both object and subject} --- as a batsman you see the bowler (object) and also examine yourself (subject) in the role-taking process. This is taken from \textbf{Charles Horton Cooley} (``I am not what I think I am\dots\ I am what I think you think I am'').
\end{itemize}
\end{KeyConcept}

\begin{CriticalPoint}
\textbf{Mead ignores structures:}
\begin{itemize}
\item Mead (micro) does not give importance to structures/constraints. On a classroom, \textbf{Durkheim} would say we act as teacher/student because of the coercive \textbf{social fact}; Mead says the classroom \textit{comes into existence only because of role-taking} --- it is the people who construct the classroom, not the classroom constraining them.
\end{itemize}
\end{CriticalPoint}

\begin{CriticalPoint}
\textbf{Mead is not a positivist:}
\begin{itemize}
\item Giving stages does \textbf{not} make one a positivist. Positivism talks about the evolution of \textit{societies} and laws driving society; Mead talks about the evolution of the \textit{self} and constantly focuses on actors (positivism ignores actors). So Mead is a \textbf{non-positivist}.
\end{itemize}
\end{CriticalPoint}

\begin{QuickRevision}
\begin{itemize}
\item Symbolic interactionism = Blumer's label (Mead's idea); reading symbols.
\item Self and society twin-born; self is object + subject (Cooley's looking-glass self).
\item Mead = micro, ignores structures (vs Durkheim's social fact); non-positivist (stages $\neq$ positivism).
\end{itemize}
\end{QuickRevision}

\sectiondivider

\subsection{Status Inconsistency (Gerhard Lenski)}

\begin{KeyConcept}
\begin{itemize}
\item \textbf{Status inconsistency} (coined by \textbf{Gerhard Lenski}, used by Merton): a mismatch between different statuses. A \textbf{crorepati auto-driver} has wealth but no status; \textbf{Ranbir Kapoor} is a member of the Kapoor family but (at a low point) a flop actor --- family name/wealth vs achievement mismatch.
\end{itemize}
\end{KeyConcept}

\begin{QuickRevision}
\begin{itemize}
\item Status inconsistency (Lenski) = mismatch of statuses (crorepati auto-driver; Kapoor-but-flop).
\end{itemize}
\end{QuickRevision}

\sectiondivider

\subsection{Revision: Family Through the Six Scholars}

\begin{DataFact}
\textbf{How each scholar reads the family:}
\begin{itemize}
\item \textbf{Durkheim} --- family is a \textbf{social fact} (external and coercive of its members); its functions are to \textbf{socialise the young into norms/values} and to \textbf{integrate members}. Conflict is only temporary (like anomie). With industrialisation the joint family weakened and re-integrated as the \textbf{nuclear family} (through interdependence) --- family survives (adaptive capacity).
\item \textbf{Marx} --- family is a \textbf{form of consciousness in the superstructure} that justifies the economic base: it is the \textbf{guardian of private property} (passed to the son), socialises people into efficient workers/businessmen, and so maintains private property and produces the working class.
\item \textbf{Weber} --- family is a structure presided over by \textbf{traditional authority} (members obey by tradition), and family is an \textbf{ideal type} --- a model; in reality there are gay, lesbian, single-parent, \textbf{transgender} (Kath Weston, \textit{Families We Choose}) and \textbf{weekend} families. The fate of our times is rationalisation and \textbf{disenchantment} (concerned with self, not family).
\item \textbf{Parsons} --- family is in the \textbf{cultural system} (motivate, sustain motivation, diffuse tension, \textbf{primary socialisation}); the extended family became the nuclear family (both ideal types) --- the family is \textbf{specialised, not declined} (an ``industrially fit'' family), running on \textbf{pattern variable A} (mother expressive/caregiver, father instrumental).
\item \textbf{Merton} --- Merton never analysed the family as such; use \textbf{reference group / role model} (born a Kapoor, you orient to Bollywood; once successful, the fraternity becomes your new in-group --- out-group $\rightarrow$ in-group = social mobility); flop films $\rightarrow$ \textbf{marginality}; and \textbf{status inconsistency}.
\item \textbf{Mead} --- self and family are \textbf{twin-born}: the self is made in the context of the family, and the family is made of selves (both are one).
\end{itemize}
\end{DataFact}

\begin{QuickRevision}
\begin{itemize}
\item Durkheim (social fact; socialise + integrate; joint$\rightarrow$nuclear), Marx (superstructure; guardian of private property), Weber (traditional authority + ideal type; disenchantment), Parsons (cultural system; specialised not declined; pattern A), Merton (reference group/marginality), Mead (self \& family twin-born).
\end{itemize}
\end{QuickRevision}

\sectiondivider

\subsection{Revision: Migration Through the Three Classical Scholars}

\begin{DataFact}
\begin{itemize}
\item \textbf{Durkheim} --- migration $\rightarrow$ material density $\uparrow$ $\rightarrow$ competition $\uparrow$ $\rightarrow$ \textbf{specialised division of labour} $\rightarrow$ cooperation restored.
\item \textbf{Marx} --- migration of the \textbf{working class} seeking employment gives industrialists an \textbf{industrial reserve army} (cheap labour); migration of the bourgeoisie is celebrated, of the proletariat looked down upon; migration (and \textbf{nationalism}, a bourgeois ideology) prevents \textbf{true class consciousness}.
\item \textbf{Weber} --- pre-industrial migration was \textbf{traditional social action} (rural-to-rural, slash-and-burn agriculture, hunting-gathering); today migration is driven by \textbf{formal rationality} (mostly \textbf{distressed} migration --- pushed, not pulled); \textbf{reverse migration} shows the family as a \textbf{safety valve/cushion} (one returns for agriculture, instrumentally --- the family ``revives'' for economic reasons).
\end{itemize}
\end{DataFact}

\begin{QuickRevision}
\begin{itemize}
\item Durkheim (density$\rightarrow$competition$\rightarrow$division of labour), Marx (reserve army; nationalism blocks class consciousness), Weber (traditional$\rightarrow$formal-rational; distressed migration; family as safety valve).
\end{itemize}
\end{QuickRevision}

\sectiondivider

\subsection{Weber as Bridge \& Asian Cosmopolitanism}

\begin{KeyConcept}
\textbf{Weber as the bridge:}
\begin{itemize}
\item Durkheim and Marx did \textbf{macro} analysis; Mead did \textbf{micro}; \textbf{Weber is the bridge} between the two extremes (social action $\rightarrow$ social structures). Parsons likewise moved from the structure of social action to the social system.
\end{itemize}
\end{KeyConcept}

\begin{KeyConcept}
\textbf{Cosmopolitanism and the ``thali'':}
\begin{itemize}
\item \textbf{Cosmopolitanism} (a model) = \textbf{Vasudhaiva Kutumbakam} (the world is one family). In reality, India is neither a \textbf{melting pot} (cultures blend into one --- ``mixed veg'') nor a \textbf{salad bowl} (cultures kept distinctly separate); it is a \textbf{thali} --- Hindus and Muslims stay together out of mutual respect \textit{and} mutual distance. This is \textbf{Asian cosmopolitanism} (studied by \textbf{Ashis Nandy}).
\end{itemize}
\end{KeyConcept}

\begin{QuickRevision}
\begin{itemize}
\item Weber = bridge (micro $\leftrightarrow$ macro); Parsons also (social action $\rightarrow$ system).
\item India = thali (not melting pot/salad bowl) = Asian cosmopolitanism (Ashis Nandy); cosmopolitanism = Vasudhaiva Kutumbakam.
\end{itemize}
\end{QuickRevision}

\sectiondivider

\section{Sociological Theories of Religion}

\subsection{Answer-Writing: Manifest \& Latent Functions of Security of Tenure}

The opening question (20 marks, 2016): \textbf{``Analyze the manifest and latent functions of the security of the tenure of the bureaucrats in the light of Merton's theory.''}

\begin{PYQLink}
\begin{itemize}
\item 2016 UPSC, 20 marks. The trick is to write \textbf{manifest}, \textbf{latent} \textit{and} \textbf{dysfunction} (writing dysfunction gets more marks).
\end{itemize}
\end{PYQLink}

\begin{DataFact}
\textbf{The background (GS detail):}
\begin{itemize}
\item Security of tenure: under new rules a bureaucrat is secure for a minimum \textbf{2 years}. It arose from \textbf{Prakash Singh v. Union of India} (police reforms, $\sim$2006) and \textbf{T. S. R. Subramanian \& others v. Union of India} (a PIL by the cabinet secretary that established the 2-year tenure). The \textbf{Second ARC} mentions \textbf{transfer raj}; politicians and corporates own the ``transfer industry'' (transfer as a weapon to sanction bureaucrats).
\end{itemize}
\end{DataFact}

\begin{ExampleBox}
\textbf{Manifest functions (anticipated):}
\begin{itemize}
\item Fosters \textbf{professionalism, good governance and efficiency}; fosters \textbf{specialisation} (a generalist becomes a specialist after 2 years in one department); gives \textbf{functional freedom} to the bureaucracy (curbing irregular transfers); \textbf{curbs the transfer industry} (weakening the politician--corporate nexus).
\end{itemize}
\end{ExampleBox}

\begin{ExampleBox}
\textbf{Latent functions (unanticipated):}
\begin{itemize}
\item Curb \textbf{alienation} and \textbf{egoistic currents} (loss of integration; bureaucrat suicides are often political); strengthens \textbf{familial ties} (children's education is not disturbed); \textbf{restores normal division of labour} and solidarity between politicians and bureaucrats. The Supreme Court's 2013 ruling also proposed a \textbf{Civil Services Board} (not yet established) to oversee transfers --- which would, however, \textbf{legitimise surveillance} of bureaucrats.
\end{itemize}
\end{ExampleBox}

\begin{ExampleBox}
\textbf{Dysfunctions:}
\begin{itemize}
\item Develops \textbf{bureaucratic personality} (rigidity, red-tapism, lackadaisical attitude); \textbf{concentration of power}; too much proximity between politicians and bureaucrats may create a nexus. (Also: depersonalisation/disenchantment.)
\end{itemize}
\end{ExampleBox}

\begin{CriticalPoint}
\textbf{Balance GS and sociology:}
\begin{itemize}
\item GS detail (ARC, Nolan Committee, the two cases) helps, but too much GS gives a non-sociological impression. Write both, and go beyond the bureaucrat's office life (family, public life) to find latent functions.
\end{itemize}
\end{CriticalPoint}

\begin{QuickRevision}
\begin{itemize}
\item Define manifest (anticipated) \& latent (unanticipated) $\rightarrow$ background (2-year tenure, Prakash Singh, TSR Subramanian, 2nd ARC transfer raj) $\rightarrow$ manifest (professionalism, specialisation, freedom, curb transfer industry) $\rightarrow$ latent (curb alienation, family ties, restore division of labour) $\rightarrow$ dysfunction (bureaucratic personality, concentration of power).
\end{itemize}
\end{QuickRevision}

\sectiondivider

\subsection{Sociological Theories of Religion: An Overview}

\begin{KeyConcept}
\textbf{What a sociologist studies:}
\begin{itemize}
\item A sociologist is not interested in whether God exists or life after death (that is philosophy/theology). Sociologists study the \textbf{impact of religion as a structure on society} --- what is religion \textit{doing} in society? (Just as we ask what the heart/brain do in the body.)
\end{itemize}
\end{KeyConcept}

\begin{DataFact}
\textbf{The theories of religion (memorise):}
\begin{enumerate}
\item \textbf{Evolutionary} theory
\item \textbf{Functional} theory
\item \textbf{Marxist} theory (+ \textbf{Neo-Marxist})
\item \textbf{Religious market} theory (also ``rational choice'' theory)
\item \textbf{Existential security} theory
\item \textbf{Interactionist} theory
\item \textbf{Feminist} theory
\item \textbf{Post-modern} theory
\end{enumerate}
\end{DataFact}

\begin{QuickRevision}
\begin{itemize}
\item Sociologist studies religion's \textit{impact on society}, not God/afterlife.
\item Eight theories: evolutionary, functional, Marxist/neo-Marxist, market, existential security, interactionist, feminist, postmodern.
\end{itemize}
\end{QuickRevision}

\sectiondivider

\subsection{Evolutionary Theory of Religion}

\begin{KeyConcept}
\begin{itemize}
\item Evolutionary theory: the more we evolve/modernise, the more pre-modern beliefs (religion) decline and science takes over as a body of knowledge.
\end{itemize}
\end{KeyConcept}

\begin{DataFact}
\textbf{Scholars:}
\begin{itemize}
\item \textbf{Auguste Comte} (Law of Three Stages): theological $\rightarrow$ metaphysical $\rightarrow$ scientific society. Religion declines with evolution.
\item \textbf{James Frazer}: the first belief system was \textbf{magic} (magicians/tantrics were rulers), then \textbf{religion}, then \textbf{science} --- three successive bodies of knowledge.
\item \textbf{Edward Burnett Tylor (E. B. Tylor)}: evolution of religion per se --- \textbf{animism $\rightarrow$ polytheism $\rightarrow$ monotheism}. The most developed societies are monotheistic --- which implicitly labels non-monotheistic societies (e.g., India with 33 crore gods) as backward (an unanticipated/latent function of the theory).
\end{itemize}
\end{DataFact}

\begin{QuickRevision}
\begin{itemize}
\item Evolutionary: more evolution $\rightarrow$ religion declines, science rises.
\item Comte (3 stages), Frazer (magic$\rightarrow$religion$\rightarrow$science), Tylor (animism$\rightarrow$polytheism$\rightarrow$monotheism).
\end{itemize}
\end{QuickRevision}

\sectiondivider

\subsection{Functional Theory: Durkheim}

\begin{KeyConcept}
\textbf{Sacred--profane dichotomy:}
\begin{itemize}
\item In every society there is a \textbf{sacred--profane dichotomy} (universal). Religion forms around \textbf{the sacred}. Durkheim's functions of religion:
\begin{enumerate}
\item \textbf{Social cohesion} --- 100 people who collectively consider something sacred are integrated around it.
\item \textbf{Collective effervescence} --- the sacred elevates you to an excited/euphoric state.
\item \textbf{Transmits common beliefs and infuses morality} among believers.
\item \textbf{Cognitive function} --- religion gave us the ability to classify the world (good/bad, pure/impure, sacred/profane, moral/immoral). Science, though developed as a refutation of religion, traces back to religion.
\end{enumerate}
\end{itemize}
\end{KeyConcept}

\begin{CriticalPoint}
\textbf{Religion will not die, it will change form:}
\begin{itemize}
\item With industrialisation/French Revolution came individualism and weakening of religion, but Durkheim (a conservative) says religion will \textbf{change its form}, not die: ``\textbf{old gods are dead, new ones are not yet born.}''
\item Religion is nothing but the \textbf{worship of society} --- as long as society exists, religion exists. The new sacred = \textbf{constitution, law, nation} (we believe in the state; the anti-nationalist is profane). Even agnostics ``believe in not believing.'' Science itself has become a new religion/sacred.
\item So Durkheim would \textbf{not} support the \textbf{secularization thesis} (decline of religious belief) --- for him, the sacred is universal.
\end{itemize}
\end{CriticalPoint}

\begin{QuickRevision}
\begin{itemize}
\item Sacred--profane is universal; 4 functions (cohesion, effervescence, morality, cognitive).
\item Religion changes form, not dies (``old gods are dead\dots''); religion = worship of society; science = new sacred.
\end{itemize}
\end{QuickRevision}

\sectiondivider

\subsection{Functional Theory: Malinowski}

\begin{DataFact}
\textbf{Who Malinowski is:}
\begin{itemize}
\item \textbf{Malinowski} is an anthropologist who conducted \textbf{field study} in the \textbf{Trobriand Islands}; he is the \textbf{inventor of participant observation} (spending years with a community, emulating them to build rapport). (Note: the first census of India, 1871, was carried out by anthropologists.)
\end{itemize}
\end{DataFact}

\begin{KeyConcept}
\textbf{Religion deals with uncertainty:}
\begin{itemize}
\item \textbf{Science} helps understand what is \textbf{certain} (boat-making with a conical end for air resistance, fertilisers for crops); \textbf{religion} deals with the \textbf{uncertain} (will my fisherman brother return from the sea? will I clear prelims?). Religion helps you overcome the uncertainty, pacify the bereaved, and deal with death (itself uncertain).
\item (Note: \textbf{participant observation} differs from a journalist's \textbf{interview} --- a Karan Thapar or Barkha Dutt spends half an hour with a person, whereas the anthropologist spends years with a community.)
\end{itemize}
\end{KeyConcept}

\begin{QuickRevision}
\begin{itemize}
\item Malinowski: Trobriand Islands, inventor of participant observation.
\item Science = certain; religion = uncertain (death, sea, exam results).
\end{itemize}
\end{QuickRevision}

\sectiondivider

\subsection{Functional Theory: Parsons}

\begin{KeyConcept}
\begin{itemize}
\item Religion is part of the \textbf{cultural system} in the AGIL paradigm. Its function (like the cultural system): \textbf{motivate individuals, sustain motivation, and diffuse tension}. Parsons' famous line: religion is a \textbf{``tonic of self-confidence.''}
\item Religion also plays an \textbf{integrative} role (like Durkheim), and \textbf{generalises values} --- it infuses morality and socialises the young into the moral order.
\end{itemize}
\end{KeyConcept}

\begin{ExampleBox}
\textbf{Morality from religion:}
\begin{itemize}
\item The auto-driver who returns the extra 100-rupee note; stopping for a puppy; empathy for the poor; the \textbf{halwa ceremony} before the Budget; the ISRO director visiting the temple before Chandrayaan/Mangalyaan --- all show that even economic and scientific actions are governed by religion.
\end{itemize}
\end{ExampleBox}

\begin{CriticalPoint}
\textbf{Parsons' source is Weber:}
\begin{itemize}
\item Parsons' idea of economic action driven by religion comes from \textbf{Max Weber} (capitalism from religion); Parsons translated Weber from German into English.
\end{itemize}
\end{CriticalPoint}

\begin{QuickRevision}
\begin{itemize}
\item Religion = cultural system (AGIL) $\rightarrow$ motivate, sustain, diffuse tension; ``tonic of self-confidence.''
\item Also integrative + generalises values (morality); examples: halwa ceremony, ISRO temple visit; source = Weber.
\end{itemize}
\end{QuickRevision}

\sectiondivider

\subsection{Marxist Theory of Religion}

\begin{KeyConcept}
\textbf{Religion is in the superstructure:}
\begin{itemize}
\item Religion is in the \textbf{superstructure} and justifies the conditions of the \textbf{economic base}. It \textbf{legitimises the illegitimate social order}.
\end{itemize}
\end{KeyConcept}

\begin{ExampleBox}
\textbf{Five points (memorise):}
\begin{enumerate}
\item \textbf{Legitimises the illegitimate order} --- poverty is justified in the name of religion (god is testing you; you'll be rich in the next birth).
\item \textbf{Creates virtue out of suffering} --- the poor are ``harijans'' (children of God), the disabled are ``divyang'' (divine), the suffering wife is ``Lakshmi''/dharampatni. (Note: ``Harijan'' was coined by \textbf{Narsinh Mehta}, popularised by Gandhi.)
\item \textbf{Opiates the masses} --- like opium, religion nullifies consciousness and thinking capacity (hence the famous line: ``religion is the \textbf{opium of the masses}'').
\item \textbf{Prevents revolutionary consciousness / revolution}.
\item \textbf{Maintains the status quo} --- religion maintains the system in favour of the dominant section (the caste structure is the making of religion).
\end{enumerate}
\end{ExampleBox}

\begin{CriticalPoint}
\textbf{Lenin \& Ambedkar:}
\begin{itemize}
\item \textbf{Lenin} called religion ``\textbf{spiritual gin}'' (like gin, it makes you feel your suffering will end).
\item \textbf{Ambedkar}: Brahmin consciousness has ruined India; his recommendation --- to become a Brahmin you must pass an exam (merit, not birth). Also: caste operates through \textbf{interpersonal interaction} --- writing one's full surname already practises caste (\textbf{Ram Manohar Lohia} on surnames).
\end{itemize}
\end{CriticalPoint}

\begin{QuickRevision}
\begin{itemize}
\item Religion = superstructure, legitimises the illegitimate order; 5 points (legitimise, virtue-from-suffering, opiate, prevent revolution, status quo).
\item Lenin: ``spiritual gin''; Ambedkar: Brahmin by exam, not birth.
\end{itemize}
\end{QuickRevision}

\sectiondivider

\subsection{Neo-Marxist Theory: Gramsci}

\begin{KeyConcept}
\textbf{Dual character of religion:}
\begin{itemize}
\item Marx sees only the \textbf{single character} of religion (justify illegitimate). Neo-Marxists see the \textbf{dual character}: religion both \textbf{maintains status quo} and \textbf{brings social change} (helps the poor free themselves). Religion integrates \textit{and} divides.
\end{itemize}
\end{KeyConcept}

\begin{ExampleBox}
\textbf{Gandhi:}
\begin{itemize}
\item Gandhi labelled Dalits ``Harijans'' (justifying poverty --- Marxist) but \textit{also} used \textbf{non-violence, satyagraha and self-renunciation} (borrowed from Jainism and Buddhism) to mobilise the masses against the British. So religion both united (Gandhi) and divided (British divide-and-rule). These are Gandhi's \textbf{praxis} (action).
\end{itemize}
\end{ExampleBox}

\begin{DataFact}
\textbf{Gramsci's labels:}
\begin{itemize}
\item Instead of Marx's ``superstructure'' and ``infrastructure'', \textbf{Gramsci} uses \textbf{civil society} (religion, media, education, family, law) and \textbf{political society}. (The Indian sociologist \textbf{Partha Chatterjee} later used these labels.)
\item \textbf{Cultural hegemony} = when the dominant section's ideas are accepted as common sense (e.g., the Hindu--Muslim divide created by the British, accepted as natural; the lockdown WhatsApp-forward about a Muslim delivery driver). The ruling class is dominant not just by owning the means of production, but by making its ideas common sense.
\item \textbf{Examples of made-common-sense ideas:} ``English is compulsory to be successful'' (spread through education, media and family --- all instruments of civil society); \textbf{Ratan Tata} as a model (owning the means of production is not enough --- dominance needs propaganda); even the idea of \textbf{only two genders (male/female)} is a dominant idea made common sense by religion (Christianity's Adam and Eve; Taoism's yin and yang).
\item \textbf{Organic intellectuals} are those who think beyond the given and bring \textbf{counter-hegemony} (Gandhi countered British hegemony by uniting Hindus and Muslims; Ambedkar countered Brahminical hegemony). \textbf{Ambedkar}: ``the people who divided us into Hindu and Muslim you call colonial, but the people who divided the Hindus you call nationalist.''
\end{itemize}
\end{DataFact}

\begin{QuickRevision}
\begin{itemize}
\item Neo-Marxism = \textbf{dual character} (religion maintains \textit{and} changes; integrates \textit{and} divides).
\item Gramsci: civil society (religion/media/education) + political society; \textbf{cultural hegemony}; \textbf{organic intellectuals} $\rightarrow$ \textbf{counter-hegemony}.
\end{itemize}
\end{QuickRevision}

\sectiondivider

\section{Theories of Religion \& Religious Practices}

\begin{CriticalPoint}
\textbf{Answer-writing tip --- 10-mark vs 20-mark questions:}
\begin{itemize}
\item A \textbf{10-mark} question is objective and to the point: define, describe features, give examples, give one counter --- done (you do \textit{not} explain). A \textbf{20-mark} question is the \textit{expanded, analytical} version of the same --- you explain a lot. Analytical questions always come as 20-markers.
\end{itemize}
\end{CriticalPoint}

\subsection{Neo-Marxist: Otto Maduro}

\begin{KeyConcept}
\begin{itemize}
\item \textbf{Otto Maduro} (neo-Marxist, like Gramsci) says \textbf{religion is independent of the economy} --- it does not always support the bourgeoisie; it also supports the oppressed/working class. Example: the \textbf{liberation theology} movement in Latin America (the priestly class supported the marginalised); \textbf{Martin Luther} also shows a religious preacher supporting the oppressed.
\end{itemize}
\end{KeyConcept}

\begin{QuickRevision}
\begin{itemize}
\item Maduro: religion is independent of economy; supports the oppressed too (liberation theology, Martin Luther).
\end{itemize}
\end{QuickRevision}

\sectiondivider

\subsection{Feminist Theory of Religion}

\begin{KeyConcept}
\begin{itemize}
\item Feminist theory is like Marxist theory with \textbf{proletariat $\rightarrow$ women, bourgeoisie $\rightarrow$ men}: women are exploited in the name of religion, promised a better life in the hereafter (their second-class citizenship on earth is justified).
\end{itemize}
\end{KeyConcept}

\begin{DataFact}
\textbf{Scholars:}
\begin{itemize}
\item \textbf{Simone de Beauvoir}: religion has exploited women; widowhood is glorified (stay widow to go to heaven); women are inferior (Manusmriti; Eve's ``original sin''; yin vs yang; the woman who breaks the ascetic's concentration; menstrual ``pollution''). \textbf{Jean Holm} compared world religions and found women's inferior status (nuns inferior to monks; women excluded from mosques).
\item Practices justified in the name of religion: \textbf{Sati, dowry, triple talaq}.
\end{itemize}
\end{DataFact}

\begin{DataFact}
\textbf{How religion is patriarchal (the criteria):}
\begin{enumerate}
\item \textbf{Within religious organisation} --- male-dominated, women excluded from priesthood.
\item \textbf{Places of worship} --- menstruation, pregnancy, childbirth considered polluting.
\item \textbf{Sacred text} --- written and interpreted by men.
\item \textbf{Religious laws and customs} --- fewer rights for women (abortion as sin, dress code/``mobile prison'' veil, punishments for adultery, inheritance laws). (The Hindu Succession Amendment Act lets daughters inherit, but only if the father made it explicit.)
\end{enumerate}
\end{DataFact}

\begin{ExampleBox}
\textbf{Exceptional women-friendly cases:}
\begin{itemize}
\item \textbf{Sabarimala} (women restricted); but there are women-only temples: \textbf{Kamakhya temple} (Assam/Andhra), \textbf{Brahma Ji temple} (Pushkar, Rajasthan --- married men cannot enter), \textbf{Mata temple} (Bihar), \textbf{Kumari Amman temple} (Kanyakumari --- only virgins/sanyasis), \textbf{Attukal Bhagwati temple} (Kerala --- Guinness record for mass female gathering). In Tamil Nadu, minister \textbf{P. K. Sekar Babu} allowed women to head 35,000 temples. But these are exceptions; the general case is patriarchal.
\end{itemize}
\end{ExampleBox}

\begin{KeyConcept}
\textbf{Subhadra Mitra Channa's Devi--Dasi dichotomy:}
\begin{itemize}
\item \textbf{Subhadra Mitra Channa} (Indian sociologist) modelled Indian women as either \textbf{Devi} (upper caste/class, celebrated) or \textbf{Dasi} (lower caste/class, subordinate/slave --- domestic help). Indian women are situated at the \textbf{intersection of religion and caste} (applicable only to Hindu religion).
\end{itemize}
\end{KeyConcept}

\begin{CriticalPoint}
\textbf{Gramsci's challenge of modernity:}
\begin{itemize}
\item ``The challenge of modernity is not to get disillusioned, but to seek through illusions'' --- to be a critical thinker is to identify whose ideas we follow.
\end{itemize}
\end{CriticalPoint}

\begin{QuickRevision}
\begin{itemize}
\item Feminist = women exploited in the name of religion (proletariat$\rightarrow$women).
\item Scholars: Simone de Beauvoir, Jean Holm; practices: sati/dowry/triple talaq; four criteria of patriarchy.
\item Devi--Dasi dichotomy (Subhadra Mitra Channa); exceptions (Kamakhya, Brahma Ji, Kumari Amman, Attukal).
\end{itemize}
\end{QuickRevision}

\sectiondivider

\subsection{Religious Market Theory}

\begin{KeyConcept}
\textbf{Stark \& Bainbridge:}
\begin{itemize}
\item \textbf{Stark and Bainbridge} (religious market/rational choice theory): religion meets a \textbf{constant, basic human need} (help in struggle, rituals at birth/marriage/death), so religion will \textbf{never disappear}.
\item Religion provides \textbf{compensators} --- promises (rewards) that compensate for unfulfilled wishes (feed birds to clear the exam; ``two minutes to heaven'').
\item Just as you choose a phone rationally (best features at minimum cost), you choose a religious guru whose advice sounds most rational/efficient/achievable --- you think in a \textbf{market-like fashion}.
\end{itemize}
\end{KeyConcept}

\begin{ExampleBox}
\textbf{Proliferation, not decline:}
\begin{itemize}
\item Instead of declining, religion \textbf{proliferates} --- ``fast religions'', instant enlightenment, digital religion (digital aarti at Banaras ghat, digital darshan, digital donation). Religion adapts to the market and the instrumental orientation of individuals. This gives rise to \textbf{sects and cults} (new world religions).
\end{itemize}
\end{ExampleBox}

\begin{CriticalPoint}
\textbf{Perpetual cycle:}
\begin{itemize}
\item Religion does not decline with evolution; it \textbf{adapts}, in a \textbf{perpetual cycle of revival and renewal} (decline $\rightarrow$ evolve $\rightarrow$ revive $\rightarrow$ new form/avatar).
\end{itemize}
\end{CriticalPoint}

\begin{QuickRevision}
\begin{itemize}
\item Stark \& Bainbridge: religion = constant need + \textbf{compensators}; people choose rationally (market).
\item Proliferation of religions, not decline; \textbf{perpetual cycle} of revival/renewal; gives rise to sects and cults.
\end{itemize}
\end{QuickRevision}

\sectiondivider

\subsection{Existential Security Theory}

\begin{KeyConcept}
\textbf{Norris \& Inglehart:}
\begin{itemize}
\item \textbf{Pippa Norris \& Ronald Inglehart} (existential security theory) challenge the religious market theory: if religion is a constant need, why are some countries more religious than others? There are \textbf{rich societies (existentially secure)} and \textbf{poor societies (existentially insecure)}.
\item Just as the economy has \textbf{boom and bust} cycles, religion has \textbf{boom and bust}: when economies boom, religion busts; when economies bust, religion booms. Poor societies (famine, disease, homelessness, no social security, 93\% informal sector) are existentially insecure, hence more religious.
\end{itemize}
\end{KeyConcept}

\begin{DataFact}
\textbf{Observation vs sociology:}
\begin{itemize}
\item Observing that poor countries are religious and rich countries are not is mere \textbf{observation}. \textbf{Sociology} (the latent/hidden part) reveals: \textbf{more religious people are born than secular people every day} --- because two-thirds of the world is poor and existentially insecure, so the world's population is increasingly religious.
\end{itemize}
\end{DataFact}

\begin{CriticalPoint}
\textbf{Sociology as martial art / debunking:}
\begin{itemize}
\item \textbf{Pierre Bourdieu}: sociology is a \textbf{martial art} (it reveals the forces of domination). \textbf{Peter Berger} calls this art \textbf{debunking} --- you are ``debunked'' when you learn something you had never thought of.
\end{itemize}
\end{CriticalPoint}

\begin{QuickRevision}
\begin{itemize}
\item Norris \& Inglehart: rich (secure) vs poor (insecure) societies; economic boom/bust $\leftrightarrow$ religious bust/boom.
\item Sociology reveals: more religious people born daily (two-thirds world poor). Bourdieu (martial art), Berger (debunking).
\end{itemize}
\end{QuickRevision}

\sectiondivider

\subsection{Postmodern Theory}

\begin{KeyConcept}
\textbf{Grace Davie: ``believing without belonging'':}
\begin{itemize}
\item \textbf{Grace Davie}: religion has become \textbf{privatised} --- no longer a moral obligation, but an individual choice. \textbf{Media has disembedded religion from its institutionalised form} --- you express religiosity through a wallpaper, a ringtone, a colour-tune, TV channels (Aastha), movies (\textit{Bruce Almighty}, \textit{Oh My God}, \textit{The Bible}), WhatsApp/Facebook/Instagram. You \textbf{believe without belonging} (you don't visit the temple, but you still believe).
\item \textbf{Robert Bellah}'s \textbf{private religion / civil religion} (Americanism) agrees with this --- religion is no longer a social fact, it is a \textbf{private/individual fact}.
\end{itemize}
\end{KeyConcept}

\begin{DataFact}
\textbf{David Lyon: Jesus in Disneyland:}
\begin{itemize}
\item \textbf{David Lyon} wrote \textit{Jesus in Disneyland} (popularised by Satish Deshpande in Delhi University): until 2018 Disneyland had Jesus-themed events/rides (a boat ride with Jesus at the centre) --- religion relocated from the church to Disneyland. This is the \textbf{relocation of religion} --- religion consumed as entertainment (Netflix web series on religious gurus).
\end{itemize}
\end{DataFact}

\begin{QuickRevision}
\begin{itemize}
\item Grace Davie: ``believing without belonging''; religion privatised; media disembedded it.
\item David Lyon: \textit{Jesus in Disneyland}; \textbf{relocation of religion}; religion as entertainment.
\end{itemize}
\end{QuickRevision}

\sectiondivider

\subsection{Interactionist Theory}

\begin{KeyConcept}
\begin{itemize}
\item Interactionist theory = look at religion \textbf{from the perspective of individuals} (a primordial form of postmodernism). Dominant names: \textbf{Clifford Geertz}, \textbf{Robert Bellah}, \textbf{Peter Berger}.
\item \textbf{Geertz}: religion helps us overcome \textbf{bafflement} and \textbf{ethical paradox} (surprise $\rightarrow$ become religious; e.g., the two-year-old surviving an earthquake).
\item \textbf{Peter Berger}: religion helps build a \textbf{universe of meaning} (you define and justify your own meaning, and call it religious).
\end{itemize}
\end{KeyConcept}

\begin{QuickRevision}
\begin{itemize}
\item Interactionism = individual's perspective (proto-postmodern); Geertz (bafflement/ethical paradox), Berger (universe of meaning), Bellah (private religion).
\end{itemize}
\end{QuickRevision}

\begin{DataFact}
\textbf{Summary: status-quoist vs social-change theories:}
\begin{itemize}
\item \textbf{Status-quoist} (religion maintains the system as it is): \textbf{functionalism, Marxism and feminism} (feminism originates in Marxism --- its conflict is men vs women, not bourgeoisie vs proletariat).
\item \textbf{Promotes social change}: \textbf{Neo-Marxism and Weber} (though for Weber only \textit{Protestantism} has the power to promote change --- other religions are status-quoist; Weber is placed here because he studied Protestantism, being himself a Protestant).
\end{itemize}
\end{DataFact}

\sectiondivider

\subsection{Religious Practices: Animism}

\begin{KeyConcept}
\textbf{Animism:}
\begin{itemize}
\item \textbf{Animism} = belief in / worship of \textbf{spirits} (of the animate and inanimate). It is the most ancient form of religion (E. B. Tylor).
\end{itemize}
\end{KeyConcept}

\begin{KeyConcept}
\textbf{E. B. Tylor's account:}
\begin{itemize}
\item Animism developed because of \textbf{intellectual needs} --- specifically to explain \textbf{death and dream} (which science cannot explain). Dream = \textbf{temporary absence of the soul} (the soul wanders in uncharted regions); death = \textbf{permanent absence of the soul} (the soul never returns). From this came belief in spirits and \textbf{ancestor worship}.
\item (Note: Malinowski also said religion develops in uncertain circumstances --- you are certain about death being inevitable, but not certain \textit{who} will die.)
\end{itemize}
\end{KeyConcept}

\begin{ExampleBox}
\textbf{Modern forms of animism:}
\begin{itemize}
\item Granting \textbf{living status} to rivers (New Zealand's Whanganui River; Ganga, Yamuna); \textbf{animal rights}; \textbf{universal human rights} (everyone has a spirit; no one is inferior --- Dalits, blacks, disabled); the \textbf{Namami Gange} programme; ``\textbf{quantum animism}'' (feeling energy in objects).
\end{itemize}
\end{ExampleBox}

\begin{QuickRevision}
\begin{itemize}
\item Animism = belief in spirits (Tylor); most ancient religion; from intellectual need to explain death \& dream.
\item Dream = temporary, death = permanent absence of soul; modern forms: river rights, animal/human rights, Namami Gange.
\end{itemize}
\end{QuickRevision}

\sectiondivider

\section{Sect, Cult \& Secularization}

\subsection{Why Sects Develop}

\begin{KeyConcept}
\textbf{Sect:}
\begin{itemize}
\item A \textbf{sect} is a group that \textbf{broke away from the mainstream religion} due to differing interpretations. Every religion has sects (Christianity: Catholics, Protestants, Mormons, Lutheranism, Calvinism; Islam: Shia, Sunni, Alawite; Hinduism: Shaivism, Vaishnavism, Shaktism, Smartism; Buddhism: Mahayana, Hinayana, Vajrayana; Jainism: Digambara, Svetambara; Sikhism: Namdhari, Sehajdhari, Udasi, Khalsa).
\end{itemize}
\end{KeyConcept}

\begin{DataFact}
\textbf{Why sects develop (internal/endogenous):}
\begin{enumerate}
\item \textbf{Ideological differences between leaders} --- e.g., \textbf{Parshvanath} (wear clothes) vs \textbf{Mahavira} (no clothes) $\rightarrow$ Svetambara vs Digambara; Mahayana (Buddha humanised) vs Hinayana (Buddha only in symbols like horse/lotus/eye); the famine in Magadha $\rightarrow$ Bhadrabahu's group went south, Sthulabhadra's stayed north $\rightarrow$ fragmentation.
\item \textbf{Spiritual/relative deprivation} --- feeling a lack of spiritualism in a material world; leaders emerge to answer ``how to fulfil obligations to the Divine.''
\end{enumerate}
\end{DataFact}

\begin{DataFact}
\textbf{External/exogenous reasons:}
\begin{itemize}
\item \textbf{Modernisation / industrialisation} (decline of tradition; critical thinking; 1517 --- Martin Luther's challenge to the Catholic Church, the era of Enlightenment).
\end{itemize}
\end{DataFact}

\begin{KeyConcept}
\textbf{Weber: sect as theodicy of the disprivileged:}
\begin{itemize}
\item \textbf{Max Weber}: a sect is the \textbf{theodicy of the disprivileged} --- the underprivileged resort to sects to answer ``why do we suffer despite hard work?'' (Protestantism for German peasants; Buddhism for Vaishya/Shudra communities).
\end{itemize}
\end{KeyConcept}

\begin{QuickRevision}
\begin{itemize}
\item Sect = broke away from mainstream religion; found in every religion.
\item Internal reasons: ideological differences + spiritual/relative deprivation; external: modernisation.
\item Weber: sect = \textbf{theodicy of the disprivileged}.
\end{itemize}
\end{QuickRevision}

\sectiondivider

\subsection{Sect vs Cult}

\begin{KeyConcept}
\textbf{Cult:}
\begin{itemize}
\item A \textbf{cult} is an organisation/movement run by an individual who \textbf{claims special spiritual knowledge} and forms around that individual (e.g., Osho, Hare Krishna movement, Transcendental Meditation, Brahmakumaris, Isha Foundation/Sadhguru, Sandeep Maheshwari). A cult leader (unlike a mere leader) claims special knowledge of the world.
\end{itemize}
\end{KeyConcept}

\begin{DataFact}
\textbf{Sect vs cult (comparison table):}\vspace{4pt}
\begin{center}
\rowcolors{2}{white}{Tint}
\begin{tabularx}{\textwidth}{p{3.2cm} p{5cm} Y}
\rowcolor{AccentDark}
\thead{Attribute} & \thead{Sect} & \thead{Cult} \\
\midrule
\textbf{Membership} & Ascribed --- you are \textit{born} into it (Shia, Sunni, Vaishnava) & Achieved --- you \textit{choose} to follow (Osho, Sadhguru) \\
\textbf{Switching} & Hard to switch (Shia$\rightarrow$Sunni is not easy) & Free --- you can belong to many cults at once \\
\textbf{Nature} & Exclusionary (cannot be both Shia and Sunni) & Inclusive (multiple cults at once) \\
\textbf{Basis of division} & Interpretation (who succeeds the Prophet), not God & Choice/individual preference \\
\textbf{Leadership} & Sustained even after the leader dies (Buddhism) & Often dies with the leader (Ram Rahim); total institution, mortification of self \\
\textbf{Legitimacy (Roy Wallis)} & \textbf{Uniquely legitimate} (monopoly of truth) & \textbf{Pluralistically legitimate} (other cults can also be legitimate) \\
\bottomrule
\end{tabularx}
\end{center}
\end{DataFact}

\begin{CriticalPoint}
\textbf{Roy Wallis \& deviance:}
\begin{itemize}
\item \textbf{Roy Wallis}: sects are \textbf{deviants} (they challenge mainstream religion --- Buddhism, Protestantism), and each sect claims a \textbf{monopoly of truth} (uniquely legitimate). Cults are \textbf{pluralistically legitimate} --- cult leaders never blame each other (no competition; they run on supply--demand like a market).
\item Cults are \textbf{total institutions} (Erving Goffman's term) --- they mortify the self; sects are more liberating.
\end{itemize}
\end{CriticalPoint}

\begin{CriticalPoint}
\textbf{Emergence of cults refutes secularization:}
\begin{itemize}
\item The daily emergence of new cults is a clear rejection of the \textbf{secularization thesis} (religion declining with advancement). Instead, old religion declines and new religion is born (Stark \& Bainbridge: religion is a constant need).
\end{itemize}
\end{CriticalPoint}

\begin{QuickRevision}
\begin{itemize}
\item Cult = forms around a charismatic individual claiming special knowledge.
\item Sect: born into, exclusionary, uniquely legitimate, sustains beyond leader. Cult: chosen, inclusive, pluralistically legitimate, total institution.
\item Rise of cults refutes secularization.
\end{itemize}
\end{QuickRevision}

\sectiondivider

\subsection{Secularization \& Its Measurement}

\begin{KeyConcept}
\textbf{Secularization:}
\begin{itemize}
\item \textbf{Secularization} = a process of decline of religious beliefs, ideas and institutions (different from \textbf{secularism}). It has been going on since the Renaissance --- science is taking over religion. Weber: work used to be a ``calling'' (a way to heaven), now it is instrumental (earning money) --- a shift from value-rational to formal-rational social action.
\end{itemize}
\end{KeyConcept}

\begin{DataFact}
\textbf{How to measure secularization (Brian Wilson):}
\begin{itemize}
\item \textbf{Bryan Wilson} is the specialist scholar to quote. Three criteria:
\begin{enumerate}
\item \textbf{Declining church membership/attendance} (point-to-point measure, like inflation).
\item \textbf{Decline of the influence of the church on the state} (loss of theocracies; Taliban Afghanistan vs democracies).
\item \textbf{Decline of the individual's religious belief}.
\end{enumerate}
\end{itemize}
\end{DataFact}

\begin{CriticalPoint}
\textbf{The problems with these measures:}
\begin{itemize}
\item Membership: not going to temple $\neq$ not religious (\textbf{believing without belonging}); going to temple $\neq$ religious (may just want peace of mind). Invalid.
\item Church--state influence: India's 2014/2019 ``Hindutva politics'' shows religion still influences elections; democracy is nurtured on religious soil. Cannot measure.
\item Individual belief: religion itself is not definite; \textbf{Robert Bellah} --- America has 300 million people, hence 300 million private religions (civil religion/Americanism/nationalism).
\end{itemize}
\end{CriticalPoint}

\begin{QuickRevision}
\begin{itemize}
\item Secularization = decline of religion ($\neq$ secularism); since Renaissance; Weber (calling$\rightarrow$money, value-rational$\rightarrow$formal-rational).
\item Measure (Brian Wilson): membership, church-state influence, individual belief --- all flawed.
\end{itemize}
\end{QuickRevision}

\sectiondivider

\subsection{Critique of Secularization}

\begin{CriticalPoint}
\textbf{Four critiques:}
\begin{enumerate}
\item \textbf{Eurocentric} --- proposed by Weber (and specialists Peter Berger, Charles Taylor); it happened in Europe first, and the same measures are wrongly applied worldwide.
\item \textbf{Durkheim} --- every society has the sacred; religion cannot decline (only changes form); decline cannot be measured.
\item \textbf{Stark \& Bainbridge (market) + David Lyon (postmodern)} --- private religions persist; cults arise to meet needs $\rightarrow$ \textbf{re-enchantment} (not disenchantment); religion changes form, develops at a rapid pace.
\item \textbf{Existential security} --- more religious people are born than secular ones (two-thirds world poor).
\end{enumerate}
\begin{itemize}
\item \textbf{Peter Berger} wrote \textit{The Sacred Canopy} and, before his death, said the secularization thesis \textbf{failed} --- we are in a \textbf{re-enchanted world}. \textbf{David Lyon} coined \textbf{re-enchantment}. Stark \& Bainbridge call it the \textbf{perpetual cycle of renewal and revival}.
\end{itemize}
\end{CriticalPoint}

\begin{QuickRevision}
\begin{itemize}
\item Secularization is eurocentric; Durkheim (sacred universal); market/postmodern (re-enchantment, perpetual cycle); existential security (more religious born).
\item Berger (Sacred Canopy) admitted the thesis failed; David Lyon: re-enchantment.
\end{itemize}
\end{QuickRevision}

\sectiondivider

\subsection{Secularism: The Indian Model}

\begin{KeyConcept}
\textbf{Secularism $\neq$ secularization:}
\begin{itemize}
\item \textbf{Secularism} is a model (not defined in the constitution, because models cannot be defined; each country has its own model). The \textbf{Indian model} runs on \textbf{equal respect for all religions / equidistance} from all religions.
\end{itemize}
\end{KeyConcept}

\begin{DataFact}
\textbf{State intervention is justified when it gives effect to equality:}
\begin{itemize}
\item The state abolished \textbf{untouchability} (a Hindu custom), \textbf{Sati}, and \textbf{triple talaq} (a Muslim custom) --- interventions justified because they give effect to equality, women's empowerment, human dignity, right to life. One must see the \textbf{intention} of intervention, not just the intervention.
\item \textbf{Santhara} (Jain fasting-unto-death): the Rajasthan High Court abolished it (Article 21), but the Supreme Court restored it (Article 25 freedom of religion) --- showing the conflict between fundamental rights.
\item Examples of selective intervention: the Indian state banned \textbf{Salman Rushdie's \textit{The Satanic Verses}} (disturbed a community's sentiments) and \textbf{Ambedkar's \textit{Riddles in Hinduism}} (a critique of upper caste) --- showing the contentious, blurred boundaries of secularism.
\end{itemize}
\end{DataFact}

\begin{QuickRevision}
\begin{itemize}
\item Secularism = model (Indian model = equal respect/equidistance). Intervention justified when it gives effect to equality (untouchability, Sati, triple talaq).
\item Santhara (Article 21 vs 25); Satanic Verses \& Riddles in Hinduism bans show blurred boundaries.
\end{itemize}
\end{QuickRevision}